\subsection{Total-depth normalization and the three paths}
\label{app:total-depth-tools}

Let $G$ be a simple tripartite graph with specified nonempty
independent parts $A,B,C$, and put $N=|V(G)|$. Complete each part
to a clique by setting
\[
 \widehat G=\left(V(G),\ E(G)\cup
                 \bigcup_{X\in\{A,B,C\}}\binom{X}{2}\right).
\]
For an integer $M\ge1$, obtain
$H^\Sigma_M(G)$ from $\widehat G$ by adding disjoint independent
sets $T_A,T_B,T_C$, with
\[
 |T_X|=M,\qquad N_{H^\Sigma_M(G)}(t)=X
 \qquad(X\in\{A,B,C\},\ t\in T_X).
\]
Depths in a rooted
forest start at one, and ancestor includes the vertex itself. An
elimination forest has one node per graph vertex and makes the
endpoints of every edge comparable in the ancestor order.
Write $\totd(F)$ for the sum of vertex depths in $F$, and
$\totd(J)$ for its minimum over elimination forests of a graph $J$.

In any elimination forest $F_G$ of $\widehat G$, the vertices of a
part are pairwise comparable because they form a clique. They lie
on one ancestor path and have a unique deepest vertex $y_X$. Define
\[
 W(F_G)=\sum_{X\in\{A,B,C\}}
             \bigl(\operatorname{depth}_{F_G}(y_X)+1\bigr),
 \qquad W(G)=\min_{F_G}W(F_G).
\]
Here and below, minima over $F_G$ range over all elimination forests
of $\widehat G$.
For a choice of part $X$, let $Y,Z$ be the other two parts. Write
$\vc(J)$ for the minimum cardinality of a vertex cover of $J$, and
define
\[
 \operatorname{cost}_X(G)=\min_{S_0}
 \left(2|S_0|+\vc\bigl(G[(Y\cup Z)\setminus S_0]\bigr)\right),
\]
where $S_0\subseteq V(G)$ covers every $X$--$(Y\cup Z)$ edge.
All cover graphs below refer to $G$, before clique completion.

\begin{lemma}[Moving one added vertex]
\label{lem:added-vertex-move}
For any $X\in\{A,B,C\}$ and $u\in T_X$, every elimination forest
of $H^\Sigma_M(G)$ can be changed so that $u$ is a leaf child of
the deepest vertex in $X$, without increasing total depth or changing
the ancestor order among vertices of $G$.
\end{lemma}

\begin{proof}
Let $F$ be the given elimination forest and $y_X$ the deepest vertex
of $X$ in $F$. The edge $uy_X$ makes $u$
and $y_X$ comparable. Delete $u$, attaching its children to its
parent, or making them roots if $u$ was a root. Reinsert $u$ as a
leaf child of $y_X$, and call the resulting forest $F_{\rm new}$.
Deletion preserves the ancestor order among
the remaining vertices. Every neighbor of $u$ belongs to $X$ and
is ancestral to $y_X$, so reinsertion preserves every edge as well.

Write $d=\operatorname{depth}_F(u)$, $e=\operatorname{depth}_F(y_X)$,
and let $D_F(u)$ be the old proper descendants of $u$. Deleting $u$
reduces each of their depths by one; reinserting it as a leaf affects
no other depth. Its new depth $d'$ satisfies $d'\le d+|D_F(u)|$:
if $u$ was below $y_X$, then $d'=e+1\le d$; otherwise $d'=e$,
and the old path from just below $u$ through $y_X$ contains
$e-d$ vertices of $D_F(u)$. The single accounting identity
\[
 \totd(F_{\rm new})-\totd(F)=d'-d-|D_F(u)|\le0
\]
therefore proves the claim.
\end{proof}

\begin{restatedlemma}{Lemma~\ref{lem:np-total-depth-added-vertices}}
Every elimination forest of $H^\Sigma_M(G)$ can be transformed without
increasing total depth so that, for every $X\in\{A,B,C\}$, all
vertices of $T_X$ are leaf children of the deepest vertex in $X$.
Consequently,
\[
\totd(H^\Sigma_M(G))
 =\min_{F_G}\bigl(\totd(F_G)+MW(F_G)\bigr),\quad MW(G)
 \le\totd(H^\Sigma_M(G))
 \le MW(G)+\frac{N(N+1)}2.
\]
\end{restatedlemma}

\begin{proof}[Proof of Lemma~\ref{lem:np-total-depth-added-vertices}]
Apply Lemma~\ref{lem:added-vertex-move} to every added vertex.
Previously moved vertices remain leaf children of their respective
deepest vertices, so this gives the stated normalization.

Deleting the added leaves gives an elimination forest $F_G$ of
$\widehat G$, and their total contribution is $MW(F_G)$.
Conversely, every $F_G$ extends by these leaves to an elimination
forest of $H^\Sigma_M(G)$. This proves the identity.

The lower bound follows from $\totd(F_G)\geq0$ and
$W(F_G)\geq W(G)$. For the upper bound,
choose $F_G$ minimizing $W$. Its total depth equals $N$ plus the
number of strict ancestor pairs, so
$\totd(F_G)\leq N+\binom N2=N(N+1)/2$.
Substituting this forest in the minimization identity proves the
upper bound.
\end{proof}

Fix an elimination forest $F_G$ of $\widehat G$. For $v\in V(G)$, put
\[
\begin{gathered}
c(v)=|\{X\in\{A,B,C\}:v\text{ is ancestral to }y_X\}|,\\S_0=\{v\in V(G):c(v)=3\},\qquad
 S_1=\{v\in V(G):c(v)=2\}.
\end{gathered}
\]
These sets record the shared prefixes of the three clique paths.

\begin{lemma}[The same pair of descendants]
\label{lem:clique-path-pair}
If $S_1$ is nonempty, all its vertices are ancestral to the same two
vertices among $y_A,y_B,y_C$.
\end{lemma}

\begin{proof}
For two vertices of a rooted forest, their sets of descendants among
$y_A,y_B,y_C$ are disjoint or one contains the other. The sets for
vertices in $S_1$ have size two, so any two intersect and must be equal.
\end{proof}

When $S_1$ is nonempty, write $X$ for the part outside this common
pair and $Y,Z$ for the other parts. When $S_1$ is empty, choose any
part as $X$ and let $Y,Z$ be the other two parts.

For a choice of $X,Y,Z$ and sets $S_0,S_1$, write
$\Gamma_X(S_0,S_1)=2|S_0|+|S_1|$.
This is the cost of a feasible pair of covers. Equivalently,
$\operatorname{cost}_X(G)$ is the minimum of $\Gamma_X(S_0,S_1)$
over sets $S_0$ covering the $X$--$(Y\cup Z)$ edges and sets
$S_1\subseteq(Y\cup Z)\setminus S_0$ covering
$G[(Y\cup Z)\setminus S_0]$.

\begin{lemma}[Covers from the three paths]
\label{lem:clique-path-covers}
The set $S_0$ covers the $X$--$(Y\cup Z)$ edges of $G$, and
$S_1\subseteq Y\cup Z$ covers $G[(Y\cup Z)\setminus S_0]$.
Moreover,
\[
 W(F_G)=N+3+\Gamma_X(S_0,S_1)
 \ge N+3+\operatorname{cost}_X(G).
\]
\end{lemma}

\begin{proof}
Every vertex of a part is ancestral to its deepest vertex. A vertex
of $S_1$ has precisely $y_Y,y_Z$ as descendants among the three
deepest vertices, so it cannot belong to $X$. A vertex of
$X\setminus S_0$ is ancestral only to $y_X$, whereas a vertex of
$(Y\cup Z)\setminus S_0$ has a nonempty descendant set contained in
$\{y_Y,y_Z\}$. These sets are disjoint, so the vertices are
incomparable and cannot be adjacent. Hence $S_0$ covers the first
set of edges. Likewise, a $YZ$-edge outside $S_0$ with neither
endpoint in $S_1$ would join vertices ancestral only to $y_Y$ and
only to $y_Z$, which are incomparable. This proves the second cover
claim.

The sum of the three depths counts each vertex $v$ exactly $c(v)$
times. Every vertex belongs to its own part's path, so $c(v)\geq1$.
The definitions of $S_0,S_1$ therefore give
\[
 \sum_{X\in\{A,B,C\}}\operatorname{depth}_{F_G}(y_X)
 =\sum_{v\in V(G)}c(v)=N+2|S_0|+|S_1|.
\]
Adding the three terminal offsets gives the identity; the cover
claims give the inequality.
\end{proof}

\begin{restatedlemma}{Lemma~\ref{lem:np-total-depth-path-cost}}
For every tripartite graph $G$ with nonempty specified parts,
\[
 W(G)=N+3+\min\{\operatorname{cost}_A(G),
                 \operatorname{cost}_B(G),
                 \operatorname{cost}_C(G)\}.
\]
\end{restatedlemma}

\begin{proof}[Proof of Lemma~\ref{lem:np-total-depth-path-cost}]
The lower bound is
\[
 W(G)=\min_{F_G}W(F_G)
 \overset{\text{Lem.~\ref{lem:clique-path-covers}}}{\ge}
 N+3+\min_{X\in\{A,B,C\}}\operatorname{cost}_X(G).
\]

For the converse, choose a part $X$ attaining the minimum, and let
$Y,Z$ be the other two parts. Choose a cover $S_0$ attaining
$\operatorname{cost}_X(G)$ and a minimum-cardinality cover $S_1$
of $G[(Y\cup Z)\setminus S_0]$. Put $S_0$ on a common chain.
Below it, put the remaining $X$ vertices on one chain. On the other
branch put $S_1$ on a chain, followed by separate chains for the
remaining $Y$ and $Z$ vertices. Empty chains are omitted. Each part
lies on one ancestor path. Every cross-part edge either meets
$S_0$ or is a remaining $YZ$-edge meeting $S_1$, so this is an
elimination forest of $\widehat G$.

The path through $S_0$ to the $X$-chain and the two paths through
$S_0,S_1$ to the $Y$- and $Z$-chains contain every vertex of their
respective parts. If a final chain is empty, end its path at the
preceding shared chain. The depth of each part's deepest vertex is
at most its path length. Counting these three path lengths counts
every vertex once, every vertex of $S_0$ twice more, and every
vertex of $S_1$ once more. Hence
\[
 W(F_G)\le N+3+\Gamma_X(S_0,S_1).
\]
The last term equals $\operatorname{cost}_X(G)$, proving the upper bound.
\end{proof}

\begin{lemma}[Normalizing successive covers]
\label{lem:total-depth-cover-normalization}
Let $X,Y,Z$ be the specified parts of $G$. Suppose that $S_0$ covers
the $X$--$(Y\cup Z)$ edges and that
$S_1\subseteq(Y\cup Z)\setminus S_0$ covers
$G[(Y\cup Z)\setminus S_0]$. An independent block of vertices in
one part with identical neighborhoods can be made wholly contained
in $S_0$, wholly contained in $S_1$, or disjoint from both, without
increasing $\Gamma_X(S_0,S_1)$. The $S_1$ state is unavailable
for a block in $X$.

Moreover, if $v\in S_0$ and every neighbor of $v$ across the
$X$--$(Y\cup Z)$ cut already belongs to $S_0$, then $v$ can be
removed from $S_0$ without increasing the cost. If $v\in Y\cup Z$,
it suffices to add $v$ to $S_1$, decreasing the cost by one.
If $v\in X$, deleting it decreases the cost by two.
\end{lemma}

\begin{proof}
Give each vertex state $0$, $1$, or $2$ according as it belongs to
$S_0$, to $S_1$, or to neither; these states cost $2$, $1$, and $0$.
For a nonempty block $D$, choose a vertex of minimum state cost in
$D$, and give every vertex of $D$ its state. This cannot increase
$\Gamma_X$. It also preserves both covers: every edge from $D$
to a vertex outside $D$ now has the same endpoint states as the
existing edge from the chosen vertex to that neighbor. The first
cover requires a state-$0$ endpoint on each first-cut edge; the second
requires a state-$0$ or state-$1$ endpoint on each $YZ$-edge.
These conditions held on the representative edge, and $D$ has no
internal edges. A block in $X$ cannot acquire state $1$, because
none of its vertices had that state. An empty block needs no change.

For the last assertion, deleting $v$ preserves the first cover by
the stated condition. If $v\in X$, the residual graph is unchanged
and the cost decreases by two. Otherwise, every new residual edge
is incident with $v$, so moving $v$ to $S_1$ covers all such edges
and replaces cost two by cost one.
\end{proof}
